\begin{textsample}{NEG Dim 1 – ap – Score 1.00 – ap/ap\_em\_60.txt}  \label{ex:f1_neg_121}
ITINERÁRIOS FORMATIVOS DA ÁREA DE MATEMÁTICA A construção dos itinerários formativos do Referencial Curricular Amapaense ( RCA ), na área de Matemática, forma um conjunto de situações e atividades educativas que os estudantes podem escolher conforme seus interesses, aptidões e objetivos.
Contemplam as competências gerais da BNCC, articuladas com as competências gerais e específicas da área, com o objetivo de aprofundar, ampliar aprendizagens, consolidar a formação integral e desenvolver habilidades que permitam aos estudantes realizarem seus projetos de vida.
PRA QUÊ? - Ampliar sua visão de mundo. - Tomar decisões e agir com autonomia e responsabilidade. - Desenvolver habilidades gerais e específicas associadas a quatro eixos estruturantes : Investigação científica, Processos criativos, Mediação e intervenção sociocultural e Empreendedorismo.

% matched lemmas: 
\end{textsample}
